% Generated draft: requires mathematical and source review.
\documentclass[11pt]{article}
\usepackage[margin=25mm]{geometry}
\usepackage{fontspec}
\setmainfont{Latin Modern Roman}
\usepackage{amsmath,amssymb,mathtools,unicode-math}
\setmathfont{Latin Modern Math}
\usepackage{enumitem,needspace,xurl,hyperref}
\setlength{\parindent}{0pt}
\setlength{\parskip}{5pt}
\setlength{\emergencystretch}{3em}
\newcommand{\N}{\mathbb{N}}
\newcommand{\Z}{\mathbb{Z}}
\newcommand{\Q}{\mathbb{Q}}
\newcommand{\R}{\mathbb{R}}
\newcommand{\C}{\mathbb{C}}
\newcommand{\F}{\mathbb{F}}
\newcommand{\A}{\mathbb{A}}
\newcommand{\PP}{\mathbb{P}}
\newcommand{\E}{\mathbb{E}}
\newcommand{\eps}{\varepsilon}
\newcommand{\dd}{\,\mathrm{d}}
\begin{document}
\section*{UnsolvedMath ID: 30004817}
\textbf{Source problem code:} catalogue entry
\subsection{Definitions and mathematical statement}
Let \(q\) be even, and let \(M\) and \(M'\) be simply connected manifolds of dimension \(2q\), equipped with normal \((q-1)\)-smoothings
\[
\bar\nu\colon M\longrightarrow B,
\qquad
\bar\nu'\colon M'\longrightarrow B.
\]
Here simply connected means path connected with trivial fundamental group. The space \(B\) is the common target of the normal structures. The supplied record does not specify its normal-structure fibration, the manifold category, or the precise convention for a normal \((q-1)\)-smoothing; these must be supplied to give the hypotheses their full meaning.

Write \(E_q(M,\bar\nu)\) for the invariant called the Q-form in the source. Its algebraic definition and the data that an isomorphism must preserve are not supplied. Thus the following is the source's conjectural implication with an unresolved invariant, rather than a self-contained definition of that invariant:
\[
\left.
\begin{gathered}
(M,\bar\nu)\text{ and }(M',\bar\nu')\text{ are bordant over }B,\\
E_q(M,\bar\nu)\cong E_q(M',\bar\nu')
\end{gathered}
\right\}
\quad\Longrightarrow\quad
(M,\bar\nu)\text{ and }(M',\bar\nu')\text{ are }h\text{-cobordant over }B.
\]
Bordism over \(B\) is understood to require extension of the specified normal structures across a cobordism. An \(h\)-cobordism is a cobordism whose two boundary inclusions are homotopy equivalences; the qualification ``over \(B\)'' requires compatibility with the specified structures over \(B\). The exact compatibility convention must be checked against the source. The record does not specify whether the manifolds are closed, compact, oriented, or allowed to have boundary, nor any relative boundary conditions or lower bound on \(q\). These omissions remain unresolved hypotheses and are not silently filled in here.
\subsection{Short English statement}
The Q-form conjecture asserts that two simply connected \(2q\)-manifolds, for even \(q\), with bordant normal \((q-1)\)-smoothings into the same \(B\) and isomorphic Q-forms are \(h\)-cobordant over \(B\). The definition of the Q-form and several geometric conventions are missing from the supplied record, so this statement requires source-based mathematical review.
\subsection{Sources}
\begin{itemize}
\item Supplied local UnsolvedMath registry record, permanent ID \(30004817\), title: The Q-Form Classification Conjecture; source problem code: \texttt{catalogue entry}. Collection: Oberwolfach Reports Open Problems. No individual proposer is identified; the supplied proposed year is 2021.
\item Supplied source citation: MATRIX-MFO Tandem Workshop: Invariants and Structures in Low-Dimensional Topology (2022), Oberwolfach Reports, pp.\ 2327--2364, DOI \texttt{10.4171/owr/2021/42}. OWR report: \url{https://doi.org/10.4171/owr/2021/42}. The supplied background describes Oberwolfach Reports as records of mathematical workshops and their open-problem discussions. The report contents have not been independently consulted here.
\item Supplied literature source, labeled ``Nagy'': \url{https://www.cambridge.org/core/journals/compositio-mathematica/article/extended-surgery-theory-for-simplyconnected-4kmanifolds/85B4255F03634F86890D193F39350770}.
\item Supplied literature source, labeled ``Nagy preprint'': \url{https://arxiv.org/abs/2402.13394}.
\item Literature assessment dated exactly 2026-08-22: the supplied curation labels the problem partially solved and states that Nagy proves the conjecture when \(H_q(B)\) is torsion-free, that later extended surgery work confirms this ``free-homology case'', and that the full conjecture in the stated generality remains unresolved. This is supplied dated curation, not an independently verified literature assessment; its time-sensitive status and cited claims require independent verification. The coefficient convention for \(H_q(B)\) is not specified in the record.
\item The supplied statement-status label is \texttt{corrected\_verified}. Its accompanying verification note says that explicit formulas were dollar-delimited and the supplied mathematical content was preserved; that note does not establish mathematical correctness or resolve the missing definitions.
\end{itemize}
\end{document}
